% =====================================================================
%  第 6 章 §6.5 B 模型与形变量子化（ch06c，前半）
%  页段：PDF 319–329（原书 p.309–319）
% =====================================================================

% [原文 p.309]
\section{B 模型与形变量子化}

\takeaway{用 Schouten 括号把多重向量场构成 $L_{\infty}$ 代数 $(\Pi C,l)$：对 Calabi–Yau 流形它形式的（Bogomolov–Tian–Todorov–Barannikov–Kontsevich），由此得到 $H$ 上的 Frobenius 流形结构，并给出小平–Spencer 写像满射的另一种说法；同一框架再用于 Poisson 结构的形变量子化。}

本节叙述 §6.4 所述理论的两个应用。一个是 B 模型中 Frobenius 结构的构造，另一个是对任意 Poisson 结构的形变量子化存在性的证明。

设 $TM$ 为 $M$ 的切丛，$\Omega^kM$ 为 $TM$ 的 $k$ 次外代数。当 $M$ 是复流形时，把 $TM^{1,0}\simeq TM$ 在 $\C$ 上的 $k$ 次外积写作 $\Omega^k_{\C}M$。置 $\Omega M=\sum\Omega^kM$，$\Omega_{\C}M=\sum\Omega^k_{\C}M$。$\Omega M$、$\Omega_{\C}M$ 的元素称为\textbf{多重向量场}（polyvector field）。在 $\Omega M$、$\Omega_{\C}M$ 上定义分次李代数结构。以后把 $\partial/\partial x_i$、$\partial/\partial z_i$ 写作 $\partial_{x_i}$、$\partial_{z_i}$。

\begin{definition}{6.138}
  在基上定义 $l_2$ 为
  \begin{align*}
    &l_2\bigl(f\partial_{x_{i_1}}\wedge\cdots\wedge\partial_{x_{i_k}},\ g\partial_{x_{j_1}}\wedge\cdots\wedge\partial_{x_{j_\ell}}\bigr)\\
    &\quad=\sum_m(-1)^{m-1}f\frac{\partial g}{\partial x_{i_m}}
      \bigl(\partial_{x_{i_1}}\wedge\cdots\wedge\widehat{\partial_{x_{i_m}}}\wedge\cdots\wedge\partial_{x_{j_\ell}}\bigr)\\
    &\qquad+\sum_m(-1)^{k+m-1}g\frac{\partial f}{\partial x_{j_m}}
      \bigl(\partial_{x_{i_1}}\wedge\cdots\wedge\widehat{\partial_{x_{j_m}}}\wedge\cdots\wedge\partial_{x_{j_\ell}}\bigr),
  \end{align*}
  并作 $\R$ 线性扩张。\footnote{按惯例，$\widehat{\partial_{x_{i_j}}}$ 之类的帽子表示去掉该项。}称 $[\xi_1,\xi_2]=(-1)^{k_1+1}l_2(\xi_1,\xi_2)$ 为 \textbf{Schouten 括号}（Schouten bracket）。（与前相同，$(-1)$ 的肩上出现 $\xi_1$ 等时，它们表示 $\deg\xi_1$ 等。）

  取复坐标 $z_i$，在上式中把 $\partial_{x_i}$ 换成 $\partial_{z_i}$，就定义了 $\Omega_{\C}M$ 的元素之间的 Schouten 括号。
\end{definition}

% [Tutor Note: Schouten 括号的两项中，$\partial_{x_{i_1}}\wedge\cdots\wedge\widehat{\partial_{x_{i_m}}}\wedge\cdots\wedge\partial_{x_{j_\ell}}$ 是把第一个自变量去掉一个因子后与第二个自变量整体外积所得的紧凑写法，故尾部是 $\partial_{x_{j_\ell}}$ 而非 $\partial_{x_{i_k}}$。若拆开写，就是 $\partial_{i_1}\wedge\cdots\widehat{\partial_{i_m}}\cdots\wedge\partial_{i_k}\wedge\partial_{j_1}\wedge\cdots\wedge\partial_{j_\ell}$。]

Schouten 括号的性质可归纳如下。

\begin{proposition}{6.139}
  $l_2$、Schouten 括号不依赖坐标而确定。设 $\xi_i\in\Omega^{k_i}M$ 或 $\xi_i\in\Omega^{k_i}_{\C}M$，则下列性质成立。

  (1) $l_2(\xi_1,\xi_2)\in\Omega^{k_1+k_2-1}M$。

  (2) $l_2(\xi_1,\xi_2)=(-1)^{k_1k_2}l_2(\xi_2,\xi_1)$。

  (3)
  \[
    l_2(l_2(\xi_1,\xi_2),\xi_3)+(-1)^{(k_1+k_2)k_3}l_2(l_2(\xi_3,\xi_1),\xi_2)
    +(-1)^{(k_2+k_3)k_1}l_2(l_2(\xi_2,\xi_3),\xi_1)=0.
  \]

  (4) $l_2(\xi_1,\xi_2\wedge\xi_3)=(-1)^{k_1k_2}\xi_2\wedge l_2(\xi_1,\xi_3)+l_2(\xi_1,\xi_2)\wedge\xi_3$。

  (5) 若 $k_1=k_2=1$，则 $l_2(\xi_1,\xi_2)$ 与通常向量场的括号积一致。

  (6) 若 $k_1=1$，$k_2=0$，则 $l_2(\xi_1,\xi_2)$ 与函数 $\xi_2$ 沿向量场 $\xi_1$ 的微分一致。

  又满足 (1)～(6) 的 $l_2$ 只限于 Schouten 括号的 $(-1)^{k_1+1}$ 倍。
\end{proposition}

命题的证明容易，留给读者。

% [原文 p.310]
\begin{remark}{6.140}
  把 Schouten 括号看作分次李代数的括号积时，须注意次数已经平移过了。通常的分次李代数 $\mathfrak{g}_{k_i}$ 的定义中，对 $X_i\in\mathfrak{g}_{k_i}$ 有 $[X_1,X_2]\in\mathfrak{g}_{k_1+k_2}$，这与命题 6.139 的 (1) 不同。因此，$[\xi_1,\xi_2]=(-1)^{k_1+1}l_2(\xi_1,\xi_2)$ 这一式子的符号，与前一节说明的 $l_2$ 和（微分）李代数的括号积的关系一致。换言之，Schouten 括号在 $\Omega M$ 上给出（关于平移了一次的次数的）分次李代数结构。

  另一方面，$\Omega M$ 有外积给出的环结构，它关于（未平移的原来的）次数分次可换。Schouten 括号对外积表现为微分，但次数上 Schouten 括号平移、外积不平移才是自然的。

  然而仔细看便会发现，(1) 意味着 $l_2$ 是 $-1$ 次的，所以这里是朝与 §6.4 相反的方向平移次数。

  为与 §6.4 的记号一致，以后考虑 $\Pi C=\Pi^2\Gamma(M;\Omega M)$。（$\Pi^2$ 不影响符号，只关系到 $l_k$ 的次数。）
\end{remark}

% [Tutor Note: 这里出现两种“平移”：Schouten 括号的 (1) 把次数降 1（$l_2$ 为 $-1$ 次），而 $\Pi^2$ 又把整体平移 2 次，使 (6.53) 的 $\Pi C$ 满足 $k+\ell=*+2$。这两次平移的合成正是让 $l_2$ 在 $\Pi C$ 上成为 0 次（即普通分次李代数括号）的关键，读下文 (6.54)～(6.60) 时务必记住这一约定。]

开始构造 B 模型中的 Frobenius 结构。设 $M$ 是复流形，
\begin{equation}
  (\Pi C)^*=\bigoplus_{k+\ell=*+2}\Gamma(M;\Omega^k_{\C}M\otimes\Lambda^{0,\ell}M) \tag{6.53}
\end{equation}
当 $\xi_i\in\Gamma(M;\Omega^{k_i}_{\C}M)$、$v\in\Gamma(M;\Lambda^{0,\ell_i}M)$ 时，置
\begin{equation}
  l_1(\xi\otimes v)=(-1)^{\xi}\xi\otimes\bar\partial v, \tag{6.54}
\end{equation}
\begin{equation}
  l_2(\xi_1\otimes v_1,\xi_2\otimes v_2)=(-1)^{v_1\xi_2}[\xi_1,\xi_2]\otimes(v\wedge v'). \tag{6.55}
\end{equation}

% [Tutor Note: (6.55) 右端的 $v\wedge v'$ 即自变量中的 $v_1\wedge v_2$，作者在此用了简写；符号 $(-1)^{v_1\xi_2}$ 是 Koszul 符号，来自把 $\xi_2$ 越过 $v_1$，若漏掉它，$l_2$ 将不满足 (2) 的分次对称性。]

% [原文 p.311]
置 $l_k=0$（$k=3,\dots$）。

\begin{lemma}{6.141}
  $( \Pi C,l)$ 是 $L_{\infty}$ 代数。
\end{lemma}

证明由命题 6.139 显然。下面的定理 6.142 由 Bogomolov [46]、Tian [365]、Todorov [367] 以及 Barannikov–Kontsevich [33] 给出。

\begin{theorem*}{6.142}
  若 $M$ 是 Calabi–Yau 流形，则 $(\Pi C,l)$ 是形式的。
\end{theorem*}

\begin{proof}
  把 §6.1 中考虑的同型 $\Omega^\ell_{\C}M\otimes\Lambda^{0,m}M\simeq\Lambda^{n-\ell,m}M$ 记作 $I$。即设 $\Omega$ 是 $\Lambda^{n,0}M$ 的正则截面、也就是全纯 $n,0$ 形式，则 $I(\xi\otimes u)=\pm i_\xi\Omega\wedge u$。

  考虑 $\partial\colon\Lambda^{\ell,m}M\to\Lambda^{\ell+1,m}M$，置 $\Delta=I^{-1}\partial I$。

  由于 $\bar\partial$ 满足 $I\bar\partial=\bar\partial I$，$\bar\partial$ 关于 $\wedge$ 是微分，即 $\bar\partial(u\wedge v)=\bar\partial u\wedge v+(-1)^uu\wedge\bar\partial v$。另一方面，$\Delta$ 与 $\wedge$ 的关系更复杂，由下面的关系式给出。

  \begin{lemma}{6.143（Tian–Todorov 恒等式）}
    \begin{equation}
      l_2(u_1,u_2)=-\bigl(\Delta(u_1\wedge u_2)-\Delta u_1\wedge u_2-(-1)^{u_1}u_1\wedge\Delta u_2\bigr). \tag{6.56}
    \end{equation}
  \end{lemma}

  \begin{proof}
    对 $\Lambda^{0,\ell}$ 成分，$I,\partial$ 的性状容易看出，故为节省记号，以下设 $u_1\in\Omega^{k_1}_{\C}$，$u_2\in\Omega^{k_2}_{\C}$。

    设 $\Omega=h\,dz_1\cdots dz_n$，则 $h$ 是全纯函数。换坐标，可设在所考虑的点处 $h=1$，$dh=0$。设 $f$ 是函数，则
    \[
      l_2(f,\partial_{z_{i_1}}\cdots\partial_{z_{i_k}})
      =\sum_j(-1)^{j-1}\partial_{z_{i_j}}(f)\partial_{z_{i_1}}\cdots\widehat{\partial_{z_{i_j}}}\cdots\partial_{z_{i_k}},
    \]
    且
    \begin{align*}
      &\bigl(\Delta(f\partial_{z_{i_1}}\cdots\partial_{z_{i_k}})-\Delta f\wedge\partial_{z_{i_1}}\cdots\partial_{z_{i_k}}
      -f\wedge\Delta(\partial_{z_{i_1}}\cdots\partial_{z_{i_k}})\bigr)\\
      &\qquad=I^{-1}\partial\bigl(\pm fh\,\partial_{z_1}\cdots\widehat{\partial_{z_{i_1}}}\cdots
      \widehat{\partial_{z_{i_k}}}\cdots\partial_{z_n}\bigr),
    \end{align*}
    因此对 $u_1=f$、$u_2=\partial_{z_{i_1}}\cdots\partial_{z_{i_k}}$，(6.56) 成立。

    一般情形下设 $u_1=f\partial_{z_{i_1}}\cdots\partial_{z_{i_k}}$，$u_2=g\partial_{z_{j_1}}\cdots\partial_{z_{j_m}}$，则
    \begin{align*}
      l_2(u_1,u_2)&=f\,l_2(\partial_{z_{i_1}}\cdots\partial_{z_{i_k}},g)\wedge\partial_{z_{j_1}}\cdots\partial_{z_{j_m}}\\
      &\qquad+(-1)^kg\,\partial_{z_{i_1}}\cdots\partial_{z_{i_k}}\wedge l_2(\partial_{z_{j_1}}\cdots\partial_{z_{j_m}},f)
    \end{align*}

% [原文 p.312]
    是。另一方面，
    \begin{align*}
      &-\bigl(\Delta(u_1\wedge u_2)-\Delta u_1\wedge u_2-(-1)^{u_1}u_1\wedge\Delta u_2\bigr)\\
      &=-\bigl(\Delta(fg\partial_{z_{i_1}}\cdots\partial_{z_{j_m}})
        -\Delta(f\partial_{z_{i_1}}\cdots\partial_{z_{i_k}})\wedge g\partial_{z_{j_1}}\cdots\partial_{z_{j_m}}\\
      &\qquad-(-1)^kf\partial_{z_{i_1}}\cdots\partial_{z_{i_k}}\wedge\Delta(g\partial_{z_{j_1}}\cdots\partial_{z_{j_m}})\bigr)
    \end{align*}
    是。盯一会儿便知两者一致。
  \end{proof}

  回到定理 6.142 的证明。

  回忆调和积分论。设 $\Delta_d,\Delta_\partial,\Delta_{\bar\partial}$ 分别是关于 $d,\partial,\bar\partial$ 的 Laplace 算子。即置 $\Delta_d=-d^*d-d^*d$ 等。（与 $\Delta$ 记号相似，但完全是两回事。）此时 $\Delta_d=2\Delta_\partial=2\Delta_{\bar\partial}$。$\bigoplus\Gamma(M;\Lambda^{k,\ell})$ 上关于 $\partial,\bar\partial,d$ 的上同调全部一致，且被 $\Delta_d=2\Delta_\partial=2\Delta_{\bar\partial}$ 下为 0 的微分形式即调和形式唯一表示。又存在 Green 算子 $G_d,G_\partial,G_{\bar\partial}$，使 $u-\Delta_dG_d$ 是调和的（$G_\partial,G_{\bar\partial}$ 也同样）。Green 算子 $G_d,G_\partial,G_{\bar\partial}$ 与 $d,\partial,\bar\partial$ 中任何一个都可交换。进一步，$\partial^*$ 与 $\bar\partial$ 交换，$\bar\partial^*$ 与 $\partial$ 交换。（以上证明参见 [156] 等。）

  置
  \[
    K=\operatorname{Ker}\partial\subset\bigoplus_{k,\ell}\Gamma(M;\Lambda^{k,\ell}).
  \]
  取同型 $I\colon\bigoplus\Gamma(M;\Lambda^{k,\ell})\to\Pi C$，把 $\Pi C$ 上的 $l_1,l_2$ 搬过来，在 $\bigoplus\Gamma(M;\Lambda^{k,\ell})$（及其子集 $K$）上定义 $l_1,l_2$。置 $l_3=\cdots=0$。设 $\varphi_1\colon K\to\bigoplus\Gamma(M;\Lambda^{k,\ell})\to\Pi C$ 为包含映射与 $I^{-1}$ 的合成。置 $\varphi_2,\dots=0$。

  \begin{lemma}{6.144}
    $(K,l)$ 是 $L_{\infty}$ 代数。又 $\varphi$ 是 $K\to\Pi C$ 的 $L_{\infty}$ 写像。
  \end{lemma}

  证明是显然的。

  \begin{lemma*}{6.145}
    $\varphi_1\colon K\to\Pi C$ 在同调上诱导同构（关于 $\bar\partial=l_1$）。
  \end{lemma*}

  \begin{proof}
    调和形式含于 $K$，故 $\varphi_1$ 在同调上诱导满射。来证单射性。借 $I$ 把 $\Pi C$ 的元素与微分形式等同。设 $u=\bar\partial v$，$u,v\in\Pi C$，$\partial u=0$。由下面的引理 6.146，可写 $u=\bar\partial\partial w$。因此 $u$ 作为 $K$ 的 $\bar\partial$ 上同调元素为 0。
  \end{proof}

% [原文 p.313]
\begin{lemma*}{6.146（$\partial\bar\partial$ 引理（[68]））}
  设 $u\in\Gamma(M;\Lambda^{k,\ell})$ 满足 $\partial u=\bar\partial u=0$，且落在 $\partial$、$\bar\partial$ 中某一个的像中。则 $u$ 可表为 $u=\partial\bar\partial w$。
\end{lemma*}

\begin{proof}
  设 $u=\bar\partial v$。则 $u-\partial\partial^*G_\partial u-\partial^*\partial G_\partial u$ 是调和的（调和这一性质关于 $\partial$、$\bar\partial$ 中任一个考虑都等价）。由于 $u\in\operatorname{Im}\partial$，而 $\partial$ 的像、$\partial^*$ 的像、调和形式这三者两两正交，故 $\partial^*\partial G_\partial u$ 以及 $u-\partial\partial^*G_\partial u-\partial^*\partial G_\partial u$ 为 0。即 $u=\partial\partial^*G_\partial u$。

  同样可得 $u=\bar\partial\bar\partial^*G_{\bar\partial}u$。于是利用上面所说的可交换性，
  \[
    u=\partial\partial^*G_\partial\bar\partial\bar\partial^*G_{\bar\partial}u
    =\partial\bar\partial\bigl(\partial^*G_\partial\bar\partial^*G_{\bar\partial}u\bigr)
  \]
  成立。取 $w=\partial^*G_\partial\bar\partial^*G_{\bar\partial}u$ 即可。
\end{proof}

% [Tutor Note: 引理 6.146 是 $\partial\bar\partial$ 引理的标准用法：先在 $K=\operatorname{Ker}\partial$ 里把 $\bar\partial$ 上同调换成 $\operatorname{Im}\partial$ 中的项，再对 $u=\bar\partial v$ 迭代一次。注意“调和”关于 $\partial$ 与 $\bar\partial$ 等价，依赖凯勒条件（$\Delta_d=2\Delta_\partial=2\Delta_{\bar\partial}$）；对一般近复或辛流形这一步不成立，这正是注 6.147 所强调的差别。]

\begin{remark}{6.147}
  Deligne–Griffiths–Morgan–Sullivan [68] 证明了凯勒流形的微分形式构成的微分代数 $(\Gamma(M;\Lambda^*M),d,\wedge)$ 是形式的。其证明与定理 6.142 的证明相近。引理 6.146 正是为此而用的。对辛流形，$(\Gamma(M;\Lambda^*M),d,\wedge)$ 一般不是形式的。
\end{remark}

回到定理 6.142 的证明。置 $H=K/\operatorname{Im}\partial$。在 $H$ 上赋予所有 $l_k$ 都为 0 的 $L_{\infty}$ 代数结构。

\begin{lemma*}{6.148}
  射影 $K\to H$ 是 $L_{\infty}$ 写像。
\end{lemma*}

\begin{proof}
  由引理 6.143，若 $u_1,u_2\in K$，则 $l_2(u_1,u_2)\in\operatorname{Im}\partial$。又由引理 6.146，$\operatorname{Im}\bar\partial\cap K\subset\operatorname{Im}\partial$。也就是说，$K$ 上 $l_k$ 的像全都含于 $\operatorname{Im}\partial$。故引理成立。
\end{proof}

再由引理 6.146，射影 $K\to H$ 在同调上诱导同构。以上便证明了定理 6.142。
\end{proof}

由定理 6.142 与定理 6.136，Calabi–Yau 流形的小平–Spencer 写像是满射。这可以改述如下。

对
\[
  x=\sum x(m),\qquad x(m)\in\bigoplus_{k+\ell=2m}H^{n-k,\ell}_{\bar\partial}(M),
\]
置 $\sum x(m)\otimes u^m=b_1\in H_0(C_*\otimes\C[u,u^{-1}],l_1)$，则存在 $b_k(x)$（$k=2,\dots$），使 $b'(x)=\sum T^kb_k(x)$

% [原文 p.314]
满足
\[
  l(e^{b'(x)})=\sum T^m\Bigl(\bar\partial b_m(x)+\frac{1}{2}\sum_{k+\ell=m}l_2(b_k(x),b_\ell(x))\Bigr)=0.
\]
下面的定理的证明据 Itagaki [187]。

\begin{theorem}{6.149}
  把绝对值充分小的复数代入 $T$ 时，$b(x)=\sum T^kb_k(x)$ 是 $C^\infty$ 收敛的。
\end{theorem}

置 $b(Tx)=\sum T^kb_k(x)$，则 $b$ 是从 $\mathcal{H}^{even}_B(M)$ 的 $0$ 的近旁到 $\bigoplus_{k+\ell:\mathrm{even}}H^\ell_{\bar\partial}(M;\Omega^k_{\C}M)$ 的光滑映射。

定义 $l_{1,b(x)}\colon\Pi C\to\Pi C$ 为
\[
  l_{1,b(x)}(y)=l(e^{b(x)}y)=l_1(y)+l_2(b(x),y).
\]
此时 $l_{1,b(x)}l_{1,b(x)}=0$ 成立。

把定理 6.123 应用于作为 $L_{\infty}$ 写像的射影 $K\to H$，就定出在同调上给出逆映射的 $L_{\infty}$ 写像 $\hat\varphi$。此时有 $\varphi_k(x^k)=b_k(x)$。

置
\[
  \psi_x(y)=\sum_{k\ge1}\frac{1}{k!}\varphi_{k+1}(x^ky).
\]
简化地写，就是 $\psi_x(y)=\varphi(e^xy)$。

\begin{lemma*}{6.150}
  $\psi_x$ 诱导从 $H$ 到 $l_{1,b(x)}$ 的上同调的同构。
\end{lemma*}

\begin{proof}
  由于 $\hat\varphi(e^{b(x)}y)=e^{b(x)}\varphi(e^{b(x)}y)$，由 $\hat d\hat\varphi(e^{b(x)}y)=\varphi\hat d(e^{b(x)}y)=0$ 得 $l_{1,b(x)}\psi_x(y)=l_{1,b(x)}\varphi(e^{b(x)}y)=0$。也就是说，$\psi_x(y)$ 是循环。所诱导的上同调映射为同构一事，稍后再证。
\end{proof}

\begin{remark}{6.151}
  按方案 6.130，对由 $L_{\infty}$ 代数的同伦同值 $\hat\varphi$ 导出的写像 $\varphi_*:M(C)\to M(C')$，可以证明 $l_{1,b}$ 的上同调与 $l_{1,\varphi_*(b)}$ 的上同调同构，引理 6.150 是它的特殊情形（参见 [122]）。
\end{remark}

$l_{1,b(x)}$ 的上同调曾是扩充形变空间 $\mathcal{M}(M)$ 的切空间。因此，引理 6.150 给出 $\mathcal{M}(M)$ 的切丛的一个平凡化。这个平凡化是把 $TH$ 的标准平凡化经坐标 $x\mapsto b(x)$ 搬过来的。这一点

% [Tutor Note: $\psi_x$ 把 $H$（即 $l_1=\bar\partial$ 的上同调）送到 $l_{1,b(x)}$ 的上同调，而 $l_{1,b(x)}$ 的上同调被断言为扩充形变空间 $\mathcal{M}(M)$ 在 $b(x)$ 处的切空间，因此 $\psi_x$ 相当于“用 $b(x)$ 把切空间搬回 $H$”。下面 (6.57) 的内积正是搬回来之后才与 $x$ 无关的（命题 6.154）。]

% [原文 p.315]
由
\[
  \frac{\partial}{\partial t_i}\varphi\Bigl(\exp\Bigl(\sum x_it_i\Bigr)\Bigr)
  =\varphi\Bigl(x_i\exp\Bigl(\sum x_it_i\Bigr)\Bigr)
\]
可以看出。

其次在 $\mathcal{M}(M)$ 上定度量。只需在 $l_{1,b(x)}$ 的上同调上定内积。$\Pi C$ 上的内积由下式定义：
\begin{equation}
  \langle u_1,u_2\rangle=\int_MI(u_1\wedge u_2)\wedge\Omega. \tag{6.57}
\end{equation}

\begin{lemma*}{6.152}
  \begin{equation}
    \langle\bar\partial u_1,u_2\rangle+(-1)^{u_1}\langle u_1,\bar\partial u_2\rangle=0, \tag{6.58}
  \end{equation}
  \begin{equation}
    \langle\Delta u_1,u_2\rangle+(-1)^{u_1+1}\langle u_1,\Delta u_2\rangle=0. \tag{6.59}
  \end{equation}
\end{lemma*}

\begin{proof}
  (6.58) 因 $\bar\partial$ 与 $I$ 可交换，由 Stokes 定理显然。(6.59) 如下。两边关于 $u_1,u_2$ 是双线性的，故用 1 的分解，可在一个坐标邻域中考虑。对 $A=a_1,\dots,a_k$ 这样的多重下标，记 $\partial_{z_A}=\partial_{z_{a_1}}\cdots\partial_{z_{a_k}}$，$d\bar z_A=d\bar z_{a_1}\cdots d\bar z_{a_k}$ 等。写作 $\underline{n}=1,\dots,n$。存在 $h$ 使 $\Omega=h\,dz_{\underline{n}}$。只需对 $u_1=f\partial_{z_A}\wedge d\bar z_{A'}$，$u_2=g\partial_{z_B}\wedge d\bar z_{B'}$ 证明即可。进一步可设 $A\cup A'=\underline{n}$，$B\cup B'=\underline{n}$，$A\cap A'=\{a_i\}=\{a'_i\}$，$B\cap B'=\emptyset$。此时
  \[
    (\Delta u_1\wedge u_2+(-1)^{u_1+1}u_1\wedge\Delta u_2)\wedge\Omega
    =\pm\partial_{a_i}(hfg)\,\partial_{z_{\underline{n}}}\wedge d\bar z_{\underline{n}},
  \]
  故 (6.59) 成立。
\end{proof}

\begin{proposition}{6.153}
  $\langle l_{1,b(x)}u_1,u_2\rangle+(-1)^{u_1}\langle u_1,l_{1,b(x)}u_2\rangle=0$。
\end{proposition}

由命题 6.153 可知 (6.57) 在 $H(\Pi C;l_{1,b(x)})$ 上定出内积。进一步下列事实成立。

\begin{proposition}{6.154}
  $\langle\psi_x(y_1),\psi_x(y_2)\rangle$ 不依赖于 $x$。
\end{proposition}

由命题 6.154 可知 $\langle\cdot,\cdot\rangle$ 定出平坦的黎曼度量。进一步，由于 $\langle\cdot,\cdot\rangle$ 非退化，由命题 6.153 可知 $\psi_y$ 在上同调上诱导单射。这是保留未证的引理 6.150 的一半。剩下一半即满射性，是上同调维数的半连续性的归结。（必要时把 $x\mapsto b(x)$ 的定义域取小一些。）

% [原文 p.316]
命题 6.153、6.154 的证明要用引理 6.155。

\begin{lemma*}{6.155}
  可以取 $\varphi_k$ 使 $\operatorname{Im}\varphi_1$ 由调和形式组成，且对 $k\ge2$ 有 $\operatorname{Im}\varphi_k\subset\operatorname{Im}\partial$。
\end{lemma*}

\begin{proof}
  回到定理 6.123 的证明，重看 $\varphi_k$ 的定义。使 $\operatorname{Im}\varphi_1$ 由调和形式组成是显然可以做到的。在 $H$ 上 $l_1=0$，故 $\varphi_k$ 由下式归纳地定义：
  \begin{align*}
    &\sum_{\sigma\in\mathfrak{S}_k}\sum_{k_1+k_2=k}\frac{1}{k_1!k_2!}
    l_2\bigl(\varphi_{k_1}(x_{\sigma(1)}\cdots x_{\sigma(k_1)}),
    \varphi_{k_2}(x_{\sigma(k_1+1)}\cdots x_{\sigma(k)})\bigr)\\
    &\qquad+l_1\varphi_k(x_1\cdots x_k)=0.
  \end{align*}
  把第 1 项写作 $y$，则 $l_1(y)=0$ 与 $y\in\operatorname{Im}l_1=\operatorname{Im}\bar\partial$ 已在定理 6.123 的证明中示明。又因 $y\in K$，有 $\partial y=0$。故由引理 6.146，存在 $w$ 使 $y=l_1\partial w$。取 $\varphi_k(x_1\cdots x_k)=-l_1(w)$ 即可。
\end{proof}

% [Tutor Note: 引理 6.155 的要点是“规范地取定”同伦提升：先在 $H$ 上解 $l_1\varphi_k=-\sum l_2(\varphi,\varphi)$，再用引理 6.146 把右端写成 $\partial(\cdots)$ 以消去障碍。所取 $\varphi_k$ 只差 $\operatorname{Im}\partial$ 的项，而这一差别不影响下面 (6.57) 的内积（由 (6.59) 它与 $\Delta$ 的像正交）。]

\begin{proof}[命题 6.153 的证明]
  $\langle\bar\partial u_1,u_2\rangle+(-1)^{u_1}\langle u_1,\bar\partial u_2\rangle=0$ 就是 (6.58)。另一方面，由命题 6.139 与引理 6.143，
\[
  \resizebox{0.97\linewidth}{!}{$\begin{aligned} l_2(b(x),u_1)\wedge u_2+(-1)^{u_1}u_1\wedge l_2(b(x),u_2) &=l_2(b(x),u_1\wedge u_2)\\ &=-\Delta(b(x)\wedge u_1\wedge u_2)+b(x)\wedge\Delta(u_1\wedge u_2) \end{aligned}$}
\]
  因此
  \begin{align*}
    &\langle[b(x),u_1],u_2\rangle+(-1)^{u_1}\langle u_1,[b(x),u_2]\rangle\\
    &\qquad=\pm\int_M\partial I(b(x)\wedge u_1\wedge u_2)\wedge\Omega
    \pm\langle b(x),\Delta(u_1\wedge u_2)\rangle
  \end{align*}
  （$I$ 已在定理 6.142 的证明中定义。）第 1 项由 Stokes 定理为 0。第 2 项由 $\Delta b(x)=0$ 与 (6.59) 为 0。
\end{proof}

\begin{proof}[命题 6.154 的证明]
  由引理 6.155，
  \[
    \langle\psi_x(y_1),\psi_x(y_2)\rangle=\langle y_1+\Delta w_1,y_2+\Delta w_2\rangle
  \]
  故由 $\Delta(y_i)=0$ 与 (6.59) 得命题。
\end{proof}

设 $u_1,u_2\in T_x(H)=H$。置 $u_1\circ u_2=\psi_x(u_1)\wedge\psi_x(u_2)$。不过左端把 $l_{1,b(x)}$ 的上同调群与 $H$ 等同看待。此时

\begin{lemma*}{6.156}
  $\langle u_1\circ u_2,u_3\rangle=\langle u_1,u_2\circ u_3\rangle$。
\end{lemma*}

\begin{proof}
  两边都等于下式。

% [原文 p.317]
  \[
    \int_MI\bigl(\psi_x(u_1)\wedge\psi_x(u_2)\wedge\psi_x(u_3)\bigr)\wedge\Omega.
  \]
\end{proof}

[33] 的主定理如下。

\begin{theorem}{6.157}
  $\langle\cdot,\cdot\rangle$、$\circ$ 在 $H$ 的原点近旁定出 Frobenius 流形的结构。
\end{theorem}

证明中剩下的是定义 6.92 的 (4)，\footnote{(5) 的证明容易。}这部分证明省略。参见 [33]、[243]。

设 $x\in H^1(M;\Lambda^{0,1})$。此时 $l_{1,b(x)}=\bar\partial+[b(x),\cdot]$ 给出通常意义的复结构形变。把这个复结构记作 $J_x$。此时，$(\Pi C,l_{1,b(x)})$ 与由复结构 $J_x$ 得到的 $(\Omega M\otimes\Lambda^{0,*},\bar\partial_{J_x})$ 保持积而同构。（虽然 $\Pi C\subset TM\otimes\Lambda^{0,*}\otimes\C$，$\Omega(M,J_M)\otimes\Lambda^{0,*}\subset TM\otimes\Lambda^{0,*}\otimes\C$，但这两个像并不一致。）因此，$\langle\cdot,\cdot\rangle$、$\circ$ 与 §6.1 引理 6.8 中构成的 Frobenius 代数除常数倍外一致。$\mathrm{int}_M\colon H^n(M;\Omega^n_{\C}M)\to\C$ 有只差常数倍的任意性，它对应全纯 $n,0$ 形式 $\Omega$ 的取法。

镜像对称的重要应用如下。计算 B 模型的 Frobenius 结构的积。镜像对称主张它与 A 模型的积一致。利用这一点，决定 A 模型的积（量子上积）中量子效应的部分，即“伪全纯曲线的个数”。这一原理成功适用的典型，是 [56] 中预言、经 [145]、[218]，最终在 [146]、[233]、[44] 中证明的 $\CP^4$ 中 5 次超曲面的情形。要计算 B 模型的 Frobenius 结构的积，必须改写为周期（period）的语言。这就是 Hodge 结构的形变理论（variation of Hodge structure）。此外，还要考虑 §6.3 所述、与镜像对称对应的巨大体积极限处的点，即巨大复结构极限（large complex structure limit），并在其近旁计算 Frobenius 结构 $\circ$、$\langle\cdot,\cdot\rangle$。巨大复结构极限对应怎样的（奇异的）复结构、怎样的坐标是平坦坐标等，是应当成为 Calabi–Yau 流形 Hodge 结构形变理论中心的重要问题。Hodge 结构的形变

% [原文 p.318]
理论在 [371] 中有详细叙述。它与 B 模型 Frobenius 结构的积的关系，涉及普通复结构形变的部分在 [64] 第 5 章有解说。关于扩充的部分，最近出现了 Barannikov 的论文 [32]。

其次叙述形变量子化问题 1.10 的肯定解决。这里 $L_{\infty}$ 写像与 Schouten 括号也起重要作用。

先叙述 Schouten 括号与 Poisson 结构的关系。设 $\{\cdot,\cdot\}$ 是 $M$ 上的 Poisson 结构。用局部坐标表示，写作
\begin{equation}
  \{f,g\}=\sum\omega^{ij}\frac{\partial f}{\partial x_i}\frac{\partial g}{\partial x_j}. \tag{6.60}
\end{equation}
则 $\sum\omega^{ij}\partial_i\wedge\partial_j$ 是多重向量场。也就是说，$\sum\omega^{ij}\partial_i\wedge\partial_j\in\Gamma(M;\Omega^2TM)$。

\begin{lemma*}{6.158}
  (6.60) 定出 Poisson 结构，与 Schouten 括号 $[\sum\omega^{ij}\partial_i\wedge\partial_j,\sum\omega^{ij}\partial_i\wedge\partial_j]$ 为 0 等价。
\end{lemma*}

\begin{proof}
\[
  \resizebox{0.97\linewidth}{!}{$\displaystyle \{\{x_i,x_j\},x_k\}+\{\{x_j,x_k\},x_i\}+\{\{x_k,x_i\},x_j\} =\sum_\ell\bigl(\partial_\ell(\omega^{ij})\omega^{\ell k} +\partial_\ell(\omega^{jk})\omega^{\ell i}+\partial_\ell(\omega^{ki})\omega^{\ell j}\bigr), $}
\]
  而右端与 $[\sum\omega^{ij}\partial_i\wedge\partial_j,\sum\omega^{ij}\partial_i\wedge\partial_j]$ 的 $\partial_i\wedge\partial_j\wedge\partial_k$ 的系数除常数倍外一致。
\end{proof}

% [Tutor Note: 引理 6.158 把“Poisson 结构的 Jacobi 恒等式”翻译成“Schouten 括号的自括号为 0”，即 $[\pi,\pi]=0$（$\pi=\sum\omega^{ij}\partial_i\wedge\partial_j$）。注意 (6.60) 中 $\omega^{ij}$ 是反称的，而此处叉积的符号约定与定义 6.138 一致，故系数比较只在差一个常数倍的意义下进行。]

其次，就代数的形变，按必要程度（稍许）叙述 Gerstenhaber [138] 等人的结果。设 $R$ 是有单位元的可换环，$A$ 是 $R$ 上的代数。先定义 Hochschild 复形（Hochschild complex）$(C^*(A,A),d)$。置 $C^k(A,A)=\operatorname{Hom}_R(A^{\otimes k},A)$，\footnote{这个定义与通常的定义次数差 1。$\Pi C$ 的次数与通常的相一致。}在其上定义微分 $d$ 为

% [原文 p.319]
\begin{align*}
  (d\Psi)(a_0\otimes\cdots\otimes a_k)
  &=a_0\Psi(a_1\otimes\cdots\otimes a_k)
   -\sum_{i=0}^{k-1}\Psi(a_0\otimes\cdots\otimes a_ia_{i+1}\otimes\cdots\otimes a_k)\\
  &\qquad-(-1)^{k-1}\Psi(a_0\otimes\cdots\otimes a_{k-1})a_k.
\end{align*}
$dd=0$ 的证明留给读者。Gerstenhaber 在 $(\Pi C^*(A,A),d)$ 上定出微分李代数结构。设 $\Pi$ 为 §6.3 引入的次数平移。即 $(\Pi C)^k(A,A)=\operatorname{Hom}_R(A^{k+1},A)$。

设 $\Psi_i\in(\Pi C)^{k_i}(A,A)=\operatorname{Hom}_R(A^{\otimes(k_i+1)},A)\sim\operatorname{Hom}_R(B_{k_i+1}\Pi A,\Pi A)$。把 $\Psi_i$ 按下列式子扩张为余微分 $\hat\Psi_i\colon B_*\Pi^{-1}A\to B_*\Pi A$：
\[
  \hat\Psi_i(a_1\otimes\cdots\otimes a_n)
  =\sum_{i=1}^n(-1)^{(i-1)k_i}a_1\otimes\cdots\otimes a_{i-1}
  \otimes\Psi_i(a_i\otimes\cdots\otimes a_{i+k_i})\otimes a_{i+k_i+1}\otimes\cdots\otimes a_n.
\]
（式中的符号来自把 $A$ 的元素都看作次数 1。）

\begin{definition}{6.159}
  称 $[\Psi_1,\Psi_2]=\Psi_1\circ\hat\Psi_2-(-1)^{k_1k_2}\Psi_2\circ\hat\Psi_1\in(\Pi C)^{k_1+k_2}(A,A)$ 为\textbf{Gerstenhaber 括号}（Gerstenhaber bracket）。
\end{definition}

置 $l_1=d$，并对 $\Psi_i\in(\Pi C)^{k_i}(A,A)$ 置
\[
  l_2(\Psi_1,\Psi_2)=(-1)^{k_1+1}[\Psi_1,\Psi_2],
\]
$l_3=\cdots=0$。

\begin{lemma*}{6.160}
  $(\Pi^2C(A,A),l)$ 是 $L_{\infty}$ 代数。
\end{lemma*}

\begin{proof}
  先确认 $[\cdot,\cdot]$ 满足 Jacobi 恒等式。由定义，
  \begin{align*}
    &\hat\Psi_1\circ\hat\Psi_2-\widehat{\Psi_1\circ\Psi_2}\\
    &=\sum_{1\le i\le i+k_1\le j\le n}(-1)^{(i-1)k_1+(j-k_1-1)k_2}a_1\otimes\cdots\\
    &\qquad\otimes\Psi_1(a_i\otimes\cdots\otimes a_{i+k_1})\otimes\cdots
     \otimes\Psi_2(a_j\otimes\cdots\otimes a_{j+k_2})\otimes\cdots\otimes a_n\\
    &\quad+\sum_{1\le j\le j+k_2\le i\le n}(-1)^{(i-1)k_1+(j-1)k_2}a_1\otimes\cdots\\
    &\qquad\otimes\Psi_2(a_j\otimes\cdots\otimes a_{j+k_2})\otimes\cdots
     \otimes\Psi_1(a_i\otimes\cdots\otimes a_{i+k_1})\otimes\cdots\otimes a_n.
  \end{align*}
  因此
  \[
    \hat\Psi_1\circ\hat\Psi_2-\widehat{\Psi_1\circ\Psi_2}
    =(-1)^{k_1k_2}\bigl(\hat\Psi_2\circ\hat\Psi_1-\widehat{\Psi_2\circ\Psi_1}\bigr).
  \]
  由此容易得到 Jacobi 律。

% [Tutor Note: 本块止于原书 p.319 页末：引理 6.160 的证明（“∘ 不结合”的括号说明）跨至 p.320，由后半块 ch06c2 接续；此处先行闭合 proof 环境，编译不受影响。]
\end{proof}

% CHUNK ch06c DONE

